\chapter{Conclusions}

This thesis set out to determine whether implementing cdt:ucum in an RDF library could be turned from a one-off effort into a general procedure. It produced a five-step methodology and validated it through four standalone implementations, in Java, TypeScript, C\#, and Python. Each leaves its host library unmodified, and each is tested against a single reference checklist so that the four can be read side by side. The procedure assumes no particular language and no particular library architecture, and the four cases show it holding across environments that differ in both.

The four cases also produced a finding that none of them shows alone. The obligations of the specification are fixed, but how far an extension can reach into a host from the outside is not, and that reach is what decides how each obligation gets met. It is set by two things: whether the host exposes its operators through a public point or evaluates them in code an extension cannot enter, and whether the host language permits intervention at runtime. The same specification was therefore satisfied through different mechanisms in each of the four hosts, and which mechanism applied followed from these two properties rather than from the obligations themselves.

\textbf{Future Work: } The clearest continuation is to apply the methodology to RDF libraries it has not yet met. RDF4J is the other major Java library; running the procedure there holds the language fixed and changes only the host architecture, isolating how much of the outcome that architecture alone decides. Sophia and rdf4cpp carry the procedure into Rust and C++, two languages absent from the four cases, where the runtime intervention that the Python implementation relied on has no obvious equivalent and the question of what reach remains is open. Each new host added this way sharpens the same map the methodology exists to draw.